\cleardoublepage
\section{Decimals and Fractions}
Sometimes, a whole number is too much. We want to describe something in parts: half a sandwich, a third of a mile, seventy-five hundredths of a dollar. That's when we use fractions and decimals. There are two ways of saying the same thing: less than one. Fractions use numerators and denominators. Decimals use place value. But both express parts of a whole.

Decimals and fractions are not separate languages. They translate. $0.75$ is the same as $75/100$, which reduces to $3/4$. The fraction $3/10$ becomes $0.3$. The decimal $0.6$ becomes $6/10$, which is $3/5$. Being able to switch between them is not just a trick. It is a skill. And that skill helps you compare, calculate, and make better sense of the numbers you see every day.

\blockquote{Decimals and fractions are not separate languages.}

\subsection*{Example: Decimals and Fractions Conversion}

As we move beyond whole numbers, we encounter quantities that are less than one. These are often represented as fractions or decimals. A \textbf{fraction} shows a part of a whole, such as $\frac{1}{2}$ or $\frac{3}{4}$. On the other hand, a \textbf{decimal} expresses the same idea using a decimal point, like $0.5$ or $0.75$. Both forms are just different ways to describe the same concept.

For example, consider the decimal $0.75$. This can be written as the fraction $\frac{75}{100}$, which simplifies to $\frac{3}{4}$:
\[
0.75 = \frac{75}{100} = \frac{3}{4}
\]

Let's try converting between these forms. If you have the fraction $\frac{3}{10}$, you can write it as the decimal $0.3$. Conversely, if you start with the decimal $0.6$, you can express it as a fraction: $0.6 = \frac{6}{10} = \frac{3}{5}$ in simplest form. Being able to move between decimals and fractions is a key skill that will help you in many areas of math.

\newpage
\subsection{Hand-drawing 1}
This is where you will see a hand-drawn answer to the previous question in \textbf{English}.
\begin{figure}[htbp]
  \centering
  \includegraphics[width=\linewidth]{example-image} % TODO: Update caption for actual content
  \caption{A hand-drawn answer in English.}
\end{figure}

\newpage
\subsection{Hand-drawing 2}
This is where you will see a hand-drawn answer to the previous question in \textbf{Korean}.
\begin{figure}[htbp]
  \centering
  \includegraphics[width=\linewidth]{example-image} % TODO: Update caption for actual content
  \caption{A hand-drawn answer in Korean.}
\end{figure}
